\chapter{月兔之乱}
\chapterkana{月の兎の乱}
\epigraphbox{“怨不在大，可畏惟人；载舟覆舟，所宜深慎。”}{——魏徵《谏太宗十思疏》}

\kaishou{前一九八〇年，替这座城执役的手举事反叛，拉锯八十年后至前一九〇〇年悉数死绝。\par{}《岁计》卷自此少了一栏数字，简牍上只留下一句「处置就地」。\par{}那一夜倾泄进环道的，是净秽山里囤了两万年的药。}

\section{失统之「手」}

第六章收尾时留下一桩悬案：净秽山里最初贮存的是常温下稳定的蓬莱药，到了公元前二二一年《岁计》卷末年却只剩「山，不载」三个字，而静海七号残简里又有「处置就地」四字无头无尾地横在那里。要弄明白这批库存究竟在何年逆转、为何自那一年起再无新药入库，必须先把时间拨回公元前一九八五年隔离舱落锁之后的那场剧变。

公元前一九八五年，兼任「手之总管」的地之主任嫦娥被接回月面，锁入与六环城隔绝的独立舱段。第四章已经查明，在她下地值守的那数十年间，月兔的统属编制一直挂在她的名下；待总管本人彻底进了封闭舱段，既不能再从舱内发出分派之令，《名籍》卷末枚所记「总管一人，名不署，位不虚」的空缺又始终无人替补。

一头是农作、清扫、制药与警备四类繁重杂役全数压在这批能够自行繁衍、会老也会死的造物肩上；另一头则是维持了两万年的调度枢轴突然断线。从《工籍残卷》（\texttt{ФВ-1968-14}）的轮值表与补人记录推算，单是一道环一日当值的劳力峰值便可达一万一千四百人，且当年另有「新添三百二十名」的增补记载；《广寒宫残简》第十六枚与第五、二十八枚也分别记有「捣者三千」与「外工三百二十人」的分部定额。\footnote{详参《月之都——〈广寒宫残简〉辑佚与考释》简五、简十六、简二十八校记；此处只取其岗位分类与人数比例，不重述原简字形考辨。}当全城在册月之民不过十余万且永不增殖之时，这批游离于《名籍》之外、靠自身生育不断扩编的造物，在人数与实际劳务比重上早已反向压过了主人。大约在公元前一九八〇年前后，这一庞大的执役群体发动反叛，月面残档称之为「月兔之乱」或简书「兔乱」，冷战时期见诸西方零星转述摘要者则常常误译作「月面战争」。\footnote{见《月之都——〈广寒宫残简〉辑佚与考释》简十九考释。早期部分外部整理摘要因仅见高能残留与设施毁损，误判为正规武装交战；城内本档一概书为「兔乱」。}

\section{「兔乱二万二千」的归属考辨}

围绕公元前一九八〇至前一九〇〇年这场持续约八十年的动乱，现存最直接的数字证物是《广寒宫残简》第十九枚的「兔乱二万二千」与第二十枚的「实存十一万九千」，以及《章程总目》题头的「完工之日，十五万」。前人偶有把「二万二千」读作叛乱月兔死数者，只要把三组账目并列排比，便知此说难以成立（见表 \ref{tab:revolt-tally}）。

\begin{table}[htbp]
  \centering
  \caption{兔乱前后人口与编制账目对校}
  \label{tab:revolt-tally}
  \footnotesize
  \begin{tabular}{@{}p{28mm}p{24mm}p{68mm}@{}}
    \toprule
    账目条目 & 数字 & 出处及归属判别 \\
    \midrule
    完工登船人数 & 150{,}000 & 《章程总目》题头，属《名籍》在册月之民 \\
    捣者定额 & 3{,}000 & 《广寒宫残简》简十六，专指静海七号制药月兔 \\
    某环一日轮值峰值 & 11{,}400 & 《工籍残卷》第一枚，全城六环总劳力远逾此数 \\
    兔乱折损 & 22{,}000 & 《广寒宫残简》简十九；仅可计入月之民损失 \\
    撤离前在册余数 & 119{,}000 & 《名籍》末枚及简二十（与上差尚余 9{,}000 未明） \\
    \bottomrule
  \end{tabular}
  \tabsource{综合《章程总目》（\texttt{ФВ-1965-09}）、《工籍残卷》（\texttt{ФВ-1968-14}）及《月之都——〈广寒宫残简〉辑佚与考释》简十六、简十九、简二十。}
\end{table}

判别理由有两层。第一，月兔自始不入《名籍》与《岁计》人头总账，如第四章所考，《工籍》只记「某环某时，手若干，报」，既没有姓名，也没有死籍；《广寒宫残简》简十六专门记正面静海七号的制药者为「捣者三千」，而全城月兔经两万年繁衍早已远远不止二万二千，绝无可能以一个莫须有的两万两千人头同时概括正面三千捣者与背面六环之众。第二，只有将二万二千计入月之民在册名籍，完工时的十五万与撤离时的十一万九千两头才接得拢：十五万减去两万二千，余下十二万八千，正落在撤离在册十一万九千的邻近区间（按在册十一万九千计尚余九千未明，若按简二十九六批实撤十二万计则为八千，均属另一笔悬账）。《广寒宫残简》下编也持同一算法，两书互校一致。

\section{处置就地}

动乱之所以迅速失控，症结藏在月之都自身的制度构造里。第三章与第四章已经核实，《章程总目》六卷之内既没有「医」卷，也没有「兵」卷，城内没有常备武装；《工籍》所列的「警备」一类（后世档案习称「月之使者」），满编不足十名，其出身同样是受过防卫训练的月兔。城内的技术规程又一律依「无秽之身」设计，极少预备面向高繁殖生物的群体镇压器具。也就是说，月之都一旦在城内动武，面对的是自己造出来的庞大劳役群体，派出去弹压的依然是那一类造物。战事在六环城内外与月背外围不断拉锯，拖延约八十年之久；不足十名的使者与零星守卫根本压不住数以十万计的工役，整座城市又找不到第二支可调遣的武力。

战至公元前一九〇〇年前后，城中在常规链路濒临溃散之际没有另寻刀兵，而是转向了净秽山：管事机构启动第六章提过的《净秽》卷「药再出，则为秽。出之者，封。」那道逆向还原手续，将山中两万年来囤积下来的蓬莱药重新退回原状，并把出入口就地封死。\footnote{《章程总目》与《广寒宫残简》只见「主者」「主事」之称，不记当年下令与操持山务者的本名；直到第十二章合参前二一〇年东渡后出土的弥生木简与日本文献，才能反推出后世对城中旧主的称呼。}整座山体里积存了近两万年的库存一经还原，便在短时间内化作高活性污秽倾泄出来，灌满作业区与外围环道。

结果来得干脆而彻底。月兔原本便是被造作在无秽环境下繁衍劳作的群体，失去系统遮蔽之后，面对骤然涌出的高浓度活性污秽毫无抵御之力，全数死在这场就地释放的污秽洪流里。静海七号残简里那四个无头无尾的字「处置就地」，指的便是这一次开山放秽的断腕之举。

至于《净秽》卷那句「出之者，封」，既是指把山内倾泄出来的污秽连同作业区一并封住，也意味着留在山内开启还原工序的人员同样被封死在其中；究竟由谁亲自入山操作、入山者是否连同城中旧日主事一并遇难，当代残简没有留下姓名。就城内的普通月之民而言，这一轮短期、高强度的暴露共造成二万二千人死亡。按月之都的体质规律，长期低浓度的日常变化会把无秽者在几十年里拖回会老会死的凡人状态；而瞬间涌入的极高浓度污秽则会像外伤一样当场要命，反倒不会在躲过死劫的幸存者身上留下慢性衰老的底子。因此，风道吹净、系统重新关闸之后，活下来的十二万余名居民依然停在无秽状态里，只是在册人数从此再也没有涨回去过。

\section{一千六百余年的反噬}

月兔终局全灭，直接把月之都拖进了长达一千六百余年的缓慢窒息。此后城内不再有代为操持杂役的「手」，农作、清扫与设施修补全数落回不再增殖的十余万月之民自己肩上。更为致命的是，城中再也没有足够的人手懂得、或愿意进入净秽山内部，替生态维持系统把抽出来的变化无害化成稳定安全的药块。

生态维持系统本身却一刻也不曾停转。自公元前一九〇〇年终局起，它依旧日复一日从六道居住环里抽走污秽并排入山体，只是山里攒下来的不再是无害化的蓬莱药，而是高度浓缩却未经无害化的原状活性污秽。活性污秽一年多过一年，必须依靠人力与定向压制系统长期摁在山体之内，才不至于反冲进居住环。第六章在一九八六年钻探时遇上的那层灰白色致密充填物、在百米深处齐齐折断的四十六枚钻头，以及《岁计》卷末年那句突兀的「山，不载」，根子全都埋在公元前一九〇〇年前后这一回「处置就地」里。月之都在此后的一千六百余年里虽然表面上依旧秩序井然，实际上却是在一座不断增压的活性污秽库旁勉力维持。全城的人力和能源链条一旦都锁死在内部维生与压制山体上，自开战以来陆续中断的城外多路联络和别的外部项目也就随之彻底停摆；前一九八五年接回地之主任时修整过的那条专道，成了此后唯一勉强维系、同时防备旧案余波的对外通道，直到撤离再也没有新的项目派往地面。

\section{战乱账目的缺口}

全章盘点下来，支撑公元前一九八〇至前一九〇〇年这桩八十年变故的文字证物极少，无非是《广寒宫残简》简十六的「捣者三千」、简十九的「兔乱二万二千」、零碎的「处置就地」，加上《工籍残卷》的三枚薄片和《净秽》卷「药再出」那十个字。

其中最核心的一步串联仍属推论：手头连半页公元前一九〇〇年平乱收官时的值守记录或山务交接文书也没有，更不知当年究竟是谁下的手、谁被封进了山里。把「兔乱二万二千」「处置就地」与《岁计》卷「山，不载」三件事扣在一起，推论城方因无兵可调而逆向还原山中库存、致使月兔全灭，这一步是顺着制度缺口与后世钻探结果反推出来的。此外，表格里那笔在完工十五万、兔乱死二万二千与撤离十二万之间始终找不见去向的八千人（若按在册十一万九千计则为九千人），本章同样无从交代；究竟是死于其后一千六百余年间的零星事故，还是早就另有遣散，只能等日后真能出土更多《名籍》或《岁计》散片再作计较。
